\documentclass{article}
\usepackage{graphicx} % Required for inserting images

\title{Reseña de AC8}
\author{Luis Angel López Peréz}
\date{\today}



\begin{document}
\maketitle
\begin{figure}
    \centering
    \includegraphics[width=0.5\linewidth]{image.png}
    \caption{Portada del Juego}
    \label{fig:placeholder}
\end{figure}

 
\section{Introducción}
A 15 dias para la salida de Ace Combat 8: Wings of Theve, según los gameplays de preeventa,es un juego que consigue emocionar al jugador en casi todos los niveles posibles, más allá de sus virtudes técnicas de la franquicia y jugables, como todos los Ace Combat el título está diseñado para despertar sentimientos genuinos: momentos de \textbf{serenidad}, catarsis, sorpresa e indignación reflexivos se combinan en una campaña que logra transmitir una historia bélica profundamente humana, algo que en la trilogia de la Playstation 2 fue una historia espectacular pero casi 6 juegos despues parece ya estar demasiado quemado.
 
\section{Historia y ambientación}
La trama se desarrolla en el año 2029 dentro del universo Strangereal, justo después de la Guerra del Faro narrada en la entrega anterior. El protagonista, Rex, según la descripción oficial es un piloto que termina heredando por accidente el legado de Jan Cope, un héroe de guerra cuya fama fue en realidad fabricada para motivar a los aliados e intimidar al enemigo. A partir de ahí, la narrativa explora tensiones políticas clásicas de la saga y, sobre todo, el desarrollo de relaciones humanas muy cercanas entre los miembros del escuadrón en turno en esta entrega siedo Joker y la tripulación del carguero Endurance.
 
\section{Duración y ritmo}
Uno de los puntos que resalta la reseña es la extensión de la campaña: treinta misiones y cerca de treinta horas de juego lo que inmediatamente volvería la entrega mas larga descoronando los 20 años de reinado como la entrega mas larga a Ace Combat 5, un incremento enorme respecto a las 19 misiones de la entrega anterior. Este tiempo adicional se aprovecha para dar espacio y trasfondo a los personajes, pero personalmente la historia es muy corta para las 30 misiones lo que va a hacer que la historia se sienta muy lenta para lo que es.
 
\section{Diseño de misiones y presentación}
Esto es definitivamente mi mayor queja, el diseño, los colores, brillos y controles son exacmeante iguales al AC7, la diferencia gráfica es mínima. También se destaca el doblaje japonés e ingles con subtítulos en español como típicamente mal traducido (aunque esto es mas culpa de Bandai que del juego).
\begin{figure}
    \centering
    \includegraphics[width=0.75\linewidth]{ac8.png}
    \caption{Imagen del Juego}
    \label{fig:placeholder}
\end{figure}
  \begin{figure}
      \centering
      \includegraphics[width=0.75\linewidth]{AC7.png}
      \caption{Graficos del anterior AC7 no distan del nuevo AC8}
      \label{fig:placeholder}
  \end{figure}
\section{Reciclaje de elementos y novedades}
Como crítica principal, señalo que el \textbf{sandbox} y las mecánicas centrales no han evolucionado prácticamente nada, los jefes como el Arsenal Bird regresa en forma de copia prácticamente igual, así como traer de regreso mecánicas viejas y persisten arquetipos de la franquicia, como las misiones de infiltración entre radares o los escuadrones enemigos. Aun así, se reconocen mejoras puntuales, como nuevas órdenes tácticas para los aliados y un sistema más flexible para equipar armas especiales.
 
\section{Aspectos técnicos}
Los requisitos mínimos publicados en Steam son: 
\begin{itemize}
    \item \textbf{SO:} Windows 11 (64 bits)
    \item \textbf{Procesador (CPU):} Intel Core i7-10700K o AMD Ryzen 5 3600X
    \item \textbf{Memoria (RAM):} 16 GB de RAM
    \item \textbf{Tarjeta gráfica (GPU):} NVIDIA GeForce RTX 2060 (6 GB) o AMD Radeon RX 6600 XT (8 GB)
    \item \textbf{DirectX:} Versión 12 (con Shader Model 6.6 o superior)
    \item \textbf{Almacenamiento:} 125 GB a 150 GB de espacio disponible en SSD (el juego requiere obligatoriamente un disco de estado sólido y una tarjeta gráfica compatible con trazado de rayos o \textit{ray tracing}).
\end{itemize}

Aunque la versión de PC recomendada es de RTX 5070 Ti, Ryzen 7 9850X3D y 32 GB de RAM.
 
\section{Conclusión}
En general, Ace Combat 8: Wings of Theve es un título \textbf{conservador} en términos de innovación jugable, siguendo la doctrina de la franquicia de apostar por la narrativa de la historia y la emoción. una entrega de una saga que deja una huella duradera gracias a sus personajes y a su tono de épica militar humana.
 
\end{document}
 
